% !TeX spellcheck = ru_RU
% !TEX root = vkr.tex
% Основной текст предоставленного обзора; библиография вынесена в related.bib.

\section{Обзор}
\label{sec:relatedworks}

\subsection{Цель обзора и отбор литературы}
\label{subsec:related_scope}

Цель обзора --- сопоставить способы адаптации сетевых детекторов,
сохранения исторических данных и калибровки порога, чтобы уточнить
задачу распределения ограниченной памяти. Рассматриваются работы
о непрерывном обучении, возвращающихся распределениях и разделении
данных между обучением и калибровкой. Корпус составлен из предварительно
отобранных статей.

Сравнение проводится по назначению памяти, способу адаптации,
смыслу калибровки и ограничиваемому ресурсу. Лимит разметки,
вместимость буфера и число сохраняемых моделей задают разные задачи.

\subsection{Изменение нормы и два источника ошибок}
\label{subsec:related_concepts}

Сетевой детектор принимает решение по признакам соединения или
потока: длительности, числу пакетов, объёму данных и характеристикам
активности. В работе интерес представляет изменение распределения
нормального трафика $P_t(X\mid Y=0)$, где $Y=0$ обозначает норму.
Оно отличается от появления новых атак и изменения их доли.
\emph{Возвратный дрейф} означает повторное появление известного
распределения после отсутствия, например $A\to B\to C\to A$.
Иллюстрацией такого режима является резервное копирование: редкая интенсивная
передача допустима, но после перерыва может вызывать тревоги.
Изменение нормы исследуется в OWAD, а возвращение распределений
--- в ARCUS и VLAD~\cite{OWAD2023,ARCUS2022,VLAD2023}.

Пусть $s_{\theta_t}(x)$ --- численная оценка аномальности наблюдения~$x$.
Большему значению соответствует более сильное отклонение от нормы
по оценке модели. Здесь $\theta_t$ --- параметры модели,
а $\tau_t$ --- порог, при превышении которого возникает тревога:
\begin{equation}
    \widehat y_t(x)=\mathbf{1}\{s_{\theta_t}(x)>\tau_t\}.
    \label{eq:related_decision}
\end{equation}
Оценкой может служить ошибка восстановления автоэнкодера или выход
классификатора. Дообучение изменяет оценочную функцию, а калибровка
порога выбирает границу решения при фиксированной модели.
После обновления модели для сохранённых калибровочных примеров
заново вычисляют оценки аномальности: значения, полученные
прежней моделью, могут отличаться от новых.

После дообучения ложные тревоги могут учащаться как из-за
ухудшения модели, так и из-за несоответствия прежнего порога
новым оценкам. Ухудшение на ранее изученных данных после
обучения на новых называют \emph{катастрофическим
забыванием}~\cite{MIR2019}. Если при этом оценки нормы и атак
сильнее перекрываются, одна перенастройка порога не улучшает
способность модели различать их. Во втором случае модель
сохраняет способность различать норму и атаки, но численные
значения оценок меняются: например, оценки нормального трафика
возрастают и начинают превышать прежний порог. Тогда изменение
порога может уменьшить число ложных тревог без дообучения
модели~\cite{LabelEfficient2026}.

\subsection{Адаптация и сохранение прежних режимов}
\label{subsec:related_network}

Общая задача адаптации к дрейфу уже решается в ADTCD:
дистилляция знаний, изменение весов наблюдений и фильтрация
выбросов направлены на быстрое обновление без загрязнения
обучающих данных~\cite{ADTCD2023}.

\textbf{OWAD} обнаруживает, объясняет и компенсирует изменения
нормального поведения~\cite{OWAD2023}. Историческая контрольная
выборка используется для сравнения распределений и адаптации;
число её примеров поддерживается постоянным. Calibrator
преобразует выходы модели для выявления сдвига. Достоинство
--- учёт изменений нормального поведения и затрат на экспертную
разметку. Ограничение состоит в зависимости от правильности
меток, которые эксперт присваивает новым наблюдениям:
ошибочное обозначение атаки как нормы может привести к тому,
что модель начнёт считать такое поведение допустимым.

\textbf{TAMR} сохраняет ограниченное число примеров для повторного
обучения; отбор учитывает близость к прототипу задачи, редкость
и давность~\cite{TAMR2025}. Преимущество --- сохранение
представительных и редких данных. Однако метод использует
известные метки задач, а эксперименты организованы как их
последовательное введение. Это отличается от возвращения
режимов с одной меткой \enquote{норма}.

\textbf{SSF} рассматривает как сохранение значимых, так и удаление
устаревших наблюдений~\cite{SSF2025}. Метод использует информацию
о норме и атаках, экспертную разметку и псевдометки. Его
достоинство --- явное управление актуальностью обучающей истории.
Но неактуальность режима сейчас не исключает его будущего возврата. У Guo и соавторов
буфер делится на представителей кластеров и примеры с высокой
неопределённостью~\cite{ContinualIDS2025}. Оба поднабора служат
обучению, а не независимой оценке распределения нормальных оценок.

\textbf{Mateen и ARCUS} сохраняют опыт в нескольких
автоэнкодерах~\cite{Mateen2024,ARCUS2022}. Mateen адаптирует
сетевой ансамбль с ограниченной разметкой и проверяет возврат
концепций; ARCUS выбирает модель, подходящую текущим данным,
либо создаёт новую. Модели, которые формируют близкие внутренние
представления данных, объединяются, чтобы сократить их число.
Достоинство --- повторное использование
прежнего представления нормы. В Mateen обучающая
база накапливает прежние и новые данные. ARCUS является общим
методом обнаружения аномалий, а не только сетевой IDS.

\textbf{VLAD} различает новые и возвращающиеся распределения,
организует примеры в иерархическую память и сокращает её при
превышении объёма~\cite{VLAD2023}. Сохранённые примеры повторно используются при
обновлении вариационного автоэнкодера. Это близкий предшественник
по памяти и возвратному дрейфу, но выбор порога тревоги
оставлен за пределами метода.

\textbf{CITADEL} сочетает представление маскированного автоэнкодера,
детектор новизны и ограниченную иерархическую память со
стратегическим отбором~\cite{CITADEL2025,CITADEL2026}. Сильная
сторона --- управление историей непосредственно в сетевой задаче.

\subsection{Отбор данных для хранения и повторного \texorpdfstring{\mbox{обучения}}{обучения}}
\label{subsec:related_retrieval}

В MIR буфер содержит случайно отобранную часть прошлых наблюдений.
При очередном дообучении используется часть примеров из этого
буфера: метод выбирает те, на которых ожидается наибольший рост
ошибки после обучения на новых данных~\cite{MIR2019}.
Для этого выполняется пробное обновление и сравниваются значения
функции потерь до и после него. Такой отбор помогает сохранить
качество модели на примерах, которые могут сильнее пострадать
от обновления, но требует дополнительных вычислений.
Критерий роста потерь не равнозначен риску тревоги:
сильное увеличение ошибки может оставить нормальный пример ниже
порога, а небольшое --- привести к его пересечению.
Таким образом, MIR отбирает примеры по изменению функции потерь,
а не непосредственно по появлению ложной тревоги.

CBRS защищает редкие размеченные классы от вытеснения
частыми~\cite{CBRS2020}. Его достоинства --- простота
и небольшие накладные расходы. Алгоритм распределяет память
по меткам классов. Если повседневная работа и резервное
копирование имеют общую метку \enquote{норма}, CBRS не различает
их как отдельные группы. Поэтому балансировка между нормой
и атаками сама по себе не защищает редкий нормальный режим
от вытеснения. В этих работах исследуются отбор и повторное
использование обучающих примеров, а распределение памяти
между обучением и калибровкой не рассматривается.

\subsection{Совместная адаптация модели и калибровки}
\label{subsec:related_joint}

\textbf{FADES} использует конформные сигналы дрейфа и при изменении
потока переобучает классификатор и повторно калибрует оценку
надёжности по общей истории фиксированного скользящего
буфера~\cite{FADES2026}. Approx-CCE сокращает число обучаемых
моделей, а адаптивный размер блока регулирует затраты.
Преимущество --- целостная потоковая процедура; обновления
предполагают надёжные метки. Конформный порог сигнала дрейфа
имеет другое назначение, чем порог сетевой тревоги.
Возвратный дрейф оставлен дальнейшим исследованиям.

\textbf{Statsenko и Bondarenko} объединяют вариационный автоэнкодер,
энтропийный отбор памяти и динамический порог на основе
экспоненциального скользящего среднего (EMA)~\cite{EntropyThreshold2025}.
Порог отслеживается через статистики, а распределение памяти между обучением
и отдельными калибровочными примерами не исследуется.

\textbf{Jeon и Kim} выбирают адаптацию возрастающей стоимости:
перекалибровку порога, переобучение на смеси старых и новых
данных либо обучение только на новом домене~\cite{LabelEfficient2026}.
Работа разделяет полезность изменения модели и порога и показывает,
что первые поступившие примеры могут плохо представлять
последующий поток. Однако ограничивается дополнительная
разметка атак; исходная выборка и валидация сохраняются.
Авторы признают ретроспективность проверки: валидация включает
будущие данные. Кроме того, отрицательный класс отдельных
детекторов содержит другие атаки. Поэтому доля ложноположительных
решений (FPR) включает ошибки на этих атаках и не характеризует
только ложные тревоги на нормальном трафике.

\subsection{Историческая калибровка и разделение данных}
\label{subsec:related_splitting}

\textbf{Replayed Calibration (RC)} сохраняет валидационные примеры
прошлых этапов и объединяет их с текущими для калибровки после
обучения~\cite{RC2024}. Достоинство --- явное рассмотрение
забывания калибровки; обучающая и калибровочная история имеют
разные назначения. Однако RC калибрует уверенность классификатора:
согласует вероятностные ответы с частотой правильных предсказаний.
Это отличается от выбора порога сетевой тревоги, а единый бюджет
двух видов памяти в работе не исследуется.

\textbf{BiC} уже разделяет ограниченные примеры старых классов
между обучением представления и настройкой линейной коррекции
выходов~\cite{BiC2019}. Сравнение долей показывает компромисс:
больше данных для коррекции означает меньше данных для сохранения
модели. Коррекция смещения к новым классам отличается от настройки
порога тревоги.

\textbf{Классификация Неймана--Пирсона} разделяет примеры класса~0
между обучением оценочной функции и выбором порога; Tong
и соавторы также предлагают адаптивный выбор доли~\cite{NP2020}.
Если класс~0 обозначает норму, задача состоит в уменьшении
пропусков атак при ограничении ложных тревог. Достоинство ---
явные статистические условия. Но анализ относится к независимым
стационарным данным; повторный выбор доли не сохраняет
автоматически доказанный контроль ошибки. Это ограничивает
перенос гарантий на дрейф и адаптивно отбираемую историю.

\textbf{Das и соавторы} оптимизируют разделение конечной выборки
между обучением регрессора и калибровкой конформного
интервала~\cite{OptimalSplit2026}. Цель --- уменьшить его длину
при заданном покрытии; эффективная доля зависит от модели
и данных. Преимущество --- теоретическая формализация компромисса,
но покрытие интервала отличается от FPR сетевого детектора.

\subsection{Сравнение ближайших работ}
\label{subsec:related_comparison}

Таблица~\ref{tab:related_memory} показывает, \emph{что именно
сохраняют} методы и \emph{для чего используют} историю.
Таблица~\ref{tab:related_calibration} разделяет разные смыслы
калибровки и указывает ограниченный ресурс. Названия методов
соответствуют описаниям выше; \enquote{не исследуется} означает
отсутствие соответствующего сравнения в рассмотренном протоколе.

\begingroup
\small
\setlength{\tabcolsep}{4pt}
\renewcommand{\arraystretch}{1.15}
\begin{longtable}{@{}>{\raggedright\arraybackslash}p{0.17\linewidth}
                        >{\raggedright\arraybackslash}p{0.37\linewidth}
                        >{\raggedright\arraybackslash}p{0.40\linewidth}@{}}
\caption{Какая история сохраняется для адаптации}
\label{tab:related_memory}\\
\toprule
\textbf{Работа} & \textbf{Что хранится} & \textbf{Для чего нужна история}\\
\midrule
\endfirsthead
\multicolumn{3}{l}{Продолжение таблицы~\thetable}\\
\toprule
\textbf{Работа} & \textbf{Что хранится} & \textbf{Для чего нужна история}\\
\midrule
\endhead
\bottomrule
\endfoot
OWAD & Контрольные примеры прежней нормы & Обнаружить изменение нормы и обновить модель\\
TAMR & Примеры известных задач в ограниченном буфере & Дообучать модель, сохраняя редкие и прежние данные\\
SSF & Представительные старые и отобранные новые примеры & Изучать изменения; удалять неактуальные данные\\
Mateen, ARCUS & Несколько моделей; в Mateen также обучающая история & Повторно использовать модель подходящего режима\\
VLAD & Примеры распределений в иерархической памяти & Распознавать возвращающийся режим и повторно обучаться на его примерах\\
CITADEL & Отобранные примеры в ограниченной иерархии & Сохранять значимые распределения при обновлении детектора\\
\end{longtable}
\endgroup

\begingroup
\small
\setlength{\tabcolsep}{4pt}
\renewcommand{\arraystretch}{1.15}
\begin{longtable}{@{}>{\raggedright\arraybackslash}p{0.17\linewidth}
                        >{\raggedright\arraybackslash}p{0.37\linewidth}
                        >{\raggedright\arraybackslash}p{0.40\linewidth}@{}}
\caption{Что калибруется и какой ресурс ограничен}
\label{tab:related_calibration}\\
\toprule
\textbf{Работа} & \textbf{Что настраивается} & \textbf{Ограничение ресурса}\\
\midrule
\endfirsthead
\multicolumn{3}{l}{Продолжение таблицы~\thetable}\\
\toprule
\textbf{Работа} & \textbf{Что настраивается} & \textbf{Ограничение ресурса}\\
\midrule
\endhead
\bottomrule
\endfoot
FADES & Конформная оценка надёжности и сигналы дрейфа & Общая скользящая история\\
Statsenko, Bondarenko & Порог сетевой тревоги по статистикам EMA & Ограничено число примеров для повторного обучения\\
Jeon, Kim & Порог тревоги; затем при необходимости модель & Лимит новой разметки\\
RC & Уверенность классификатора & Распределение единого бюджета двух видов памяти не исследуется\\
BiC & Коррекция смещения ответов к новым классам & Примеры старых классов делятся между обучением и коррекцией\\
Нейман--Пирсон & Порог с ограничением ошибки первого рода & Конечная выборка делится между обучением и порогом\\
Das и соавторы & Конформный регрессионный интервал & Конечная выборка делится между обучением и калибровкой\\
\end{longtable}
\endgroup

\clearpage
\subsection{Выводы}
\label{subsec:related_conclusions}

Рассмотренные работы уже исследуют replay, сохранение возвращающихся
распределений, историческую калибровку и выбор долей данных.
OWAD, TAMR и VLAD показывают, как история помогает адаптировать
модель и сохранять сведения о прежних режимах. FADES и работа
Statsenko и Bondarenko объединяют адаптацию с калибровочными
механизмами. RC использует прошлые данные для калибровки
уверенности, а BiC, классификация Неймана--Пирсона и работа
Das и соавторов рассматривают разделение данных между двумя
процедурами.

Эти направления отвечают на разные вопросы: как сохранить
способность модели распознавать прежние данные и как настроить
решения после изменения её оценок. Их результаты по отдельности
не определяют, как распределять единый ограниченный объём
исторических данных между дообучением и калибровкой порога
в изменяющемся потоке. Возвращение нормальных режимов делает
сохранённую историю вновь полезной и связывает этот вопрос
с темой работы.

% Дополнительные материалы из исходного списка литературы.
\nocite{RCCode,RCReproduction}
